\subsection{Two-party Exponential Mechanism}\label{sec:dp}
Recall that one of our main motivations for this work is to instantiate
a two-party version of the exponential mechanism to achieve differential
privacy.  We observe that for many natural loss functions (i.e., when
the loss function is additive across the two parties), the exponential
mechanism on two parties is essentially equivalent to product sampling.
We explain this further with a concrete example in
Section~\ref{sec:example}.  


\subsubsection{A concrete example}\label{sec:example}
Suppose we want to choose a classifier minimizing the $L_2$ error over a
test dataset while preserving differential privacy of the labeled
examples.  
%
Suppose there are $n$ machine learning classifiers $(c_1, \ldots, c_n)$,
and a test dataset $D = (d_1, \ldots, d_{|D|})$ consists of $|D|$ rows.
Let $\ell_j \in \{0,1\}$ be the label of the $j$-th row $d_j$ of the
dataset.
%
For a machine learning classifier $c_i$, we define its $L_2$ loss
function as follows: 

\[
f^{c_i}_{\mathsf{loss}}(D) := \sum_{j \in |D|}(c_i(d_j) - \ell_j)^2/|D|.
\]

Now, consider a two-party federated setting in which the parties would
like to perform computation on the aggregation of their local datasets.
In particular, we assume party $P_1$ (resp., party $P_2$) holds dataset
$D_1$ (resp., $D_2$) with $|D_1| = |D_2|$. Let $D = D_1 || D_2$.
%We assume the parties know the size of $|D|$.  \dnote{This should not
%violate differential privacy since (1) the test subset $D$ is usually
%randomly sampled from the entire training set and typically has far
%smaller size than the entire training set, so we can view $|D|$ as a
%parameter to the algorithm and require that both parties have exactly
%$|D|/2$ number of examples in their testing set (2) if needed, $D$
%could be padded with dummy rows to make the count differentially
%private.}\anote{I think it's easier to skip this discussion and just
%set the sets to be of size $|D|/2$.}


\paragraph{DP mechanism in the central curator model.}
In our mechanism, the central curator would receive input from parties
$P_1$ and $P_2$ and choose classifier $c_i$ with a $(\epsilon, 0)$-DP
guarantee using the exponential mechanism. 

We observe that the $L_2$ loss function $f^{c_i}_{\mathsf{loss}}(D)$
over the entire dataset $D$ can be computed by each party $P_b$ {\em
first locally computing}
\[
f^{c_i}_{\mathsf{loss}}(D_b) := \sum_{j \in |D_b|}(c_i(d_{b,j}) - \ell_{b,j})^2/|D|,
\]
and then computing 
$f^{c_i}_{\mathsf{loss}}(D) = f^{c_i}_{\mathsf{loss}}(D_1) + f^{c_i}_{\mathsf{loss}}(D_2).$


Based on the above observation, in our mechanism, each party $P_b$
computes a vector $\vec{v}_b$ as follows: 
%
\BI
\item[] For $b \in \{1,2\}$, let $\vec{v}_b = (v_{b,1}, \ldots, v_{b,
n})$, where $v_{b,i} = e^{-\epsilon \cdot \frac{
f^{c_i}_{\mathsf{loss}}(D_b)}{2 \Delta u}}$ and $\Delta u =
\frac{\e}{40(\log n + t)}$, and $\secparam$ is the security parameter. 
\EI
%
Then, each party $P_b$ computes $\w_b :=
\frac{\vec{v}_{b}}{||\vec{v}_b||_1}$ (i.e., the normalization of
$\vec{v}_b$), and sends $\w_b$ to the central curator. 
%
Finally, the curator will choose classifier $c_i$ with the probability
$\frac{w_{i,i} \cdot w_{i,2}}{\ipw}$. 

\begin{lemma}
If $|D| \ge 40 (\log n + \secparam)/\epsilon$, our mechanism
provides $(\epsilon, 0)$-DP.  
\end{lemma}

\begin{proof} 
Let $q_i(D) = \frac{-f^{c_i}_{\mathsf{loss}}(D)}{2 \Delta u}$. 
We first show that drawing a sample from the product distribution of
$\vec{w}_1, \vec{w}_2$ is identical to running the exponential mechanism
to select classifier $c_i$ with probability $$\frac{e^{\epsilon q_i(D)}
}{\sum_{j \in [n]} e^{\epsilon q_j(D)}}.$$ 
Recall that product sampling on inputs $\vec{w}_1,
\vec{w}_2$ returns index $i$ with probability $\frac{{w}_{1,i} \cdot
{w}_{2,i}}{\langle \vec{w}_1, \vec{w}_2\rangle}$.  Since
${w}_{b,i} = \frac{{v}_{b,i}}{||\vec{v}_b||_1}$, we can rewrite
the previous equation as
\begin{align*} \frac{{w}_{1,i} \cdot {w}_{2,i}}{\langle \vec{w}_1,
\vec{w}_2\rangle} & = \frac{{v}_{1,i}
\cdot {v}_{2,i}}{\ip{\vec{v}_1}{\vec{v}_2}} 
    &= \frac{e^{\epsilon \cdot q_i(D_1)}\cdot e^{\epsilon \cdot q_i(D_2)}}
    {\sum_{j \in [n]} e^{\epsilon \cdot q_j(D_1)} \cdot e^{\epsilon \cdot q_j(D_2)}}
    = \frac{e^{\epsilon \cdot q_i(D)}} {\sum_{j \in [n]} e^{\epsilon \cdot q_j(D)}} 
\end{align*} 

Note that the loss functions have global sensitivity of $1/|D|$ since a
change to the $j$-th row of $D_b$ can cause
$f^{c_i}_{\mathsf{loss}}(D_b)$ to change by $\pm 1/|D|$.
%
Instead of using $1/|D|$, our mechanism uses a larger value $\Delta u$.
Since the $\Delta u$ is larger than the sensitivity whenever $|D| \ge 40 (\log{n}+\secparam)/\epsilon$, and our mechanism
is a simple exponential mechanism, $(\e, 0)$-DP holds.  
\qed
\end{proof}

% \dnote{Notes so we don't forget:
% $\frac{\vec{w}_1[i] \cdot \vec{w}_2[i]}{\langle \vec{w}_1, \vec{w}_2\rangle}$
% is equivalent to:
% $\frac{\vec{w}_1[i] \cdot \vec{w}_2[i]}{||\vec{w}_1||_1 \cdot ||\vec{w}_2||_1\langle \vec{w}_1/||\vec{w}_1||_1, \vec{w}_2/||\vec{w}_2||_1\rangle}$}

%
%We say that $f^{c_i}_\mathsf{loss}$ has
%\emph{low-sensitivity} if
%Note the global sensitivity of $f^{c_i}_\mathsf{loss}$ is at most $1/|D|$.
%\anote{I thought about generalizing this, but then
%the error below gets really messy.  Should we generalize?}
%

\paragraph{Utility.}
Although using a larger value $\Delta u$ deteriorates the utility of the
mechanism, we show that the utility is still acceptable.  Applying
Theorem~3.11
in~\cite{DR14} to our setting, we have $$\Pr\left[f_{\sf loss}^{c_i}(D)
\le f_{\sf loss}^{c_{opt}}(D) - \frac{2 \Delta u}{\e} \cdot (\log n + \secparam)
\right] \le e^{-\secparam}.$$ 

Noting that $\Delta u = \frac{\e}{40(\log
n + \secparam)}$, we have $$\Pr\left[f_{\sf loss}^{c_i}(D) \le f_{\sf
loss}^{c_{opt}}(D) - 1/20 \right] \le e^{-\secparam}.$$ 
Viewing $f_{\sf
loss}^{c_{opt}}(D)$ as the optimal accuracy for the chosen classifier.  This implies that our mechanism returns a classifier that is at most 5\% less accurate than the optimal classifier.  We note that an even smaller loss in accuracy can be achieved by increasing $\Delta u$ and the minimum size of $D$ accordingly.

Jumping ahead, we use this larger $\Delta u$ in order to achieve
differential privacy of the approximate inner product evaluation to be
described in the next section.

\subsubsection{Differentially-Private Inner Product for the Exponential Mechanism}

\paragraph{Issue: DP is broken due to leakage $\ipw$.}
Based on the result of the previous subsection, we can simply run the
product sampling protocol to achieve a two-party exponential mechanism
without the central curator. However, there is one issue we need to
address.
%
In particular, the leakage from the previously described protocols for
product sampling violates the DP guarantee of the exponential mechanism;
the leakage $\ipw$ is clearly not differentially private with respect to
$P_2$'s input.

Thus, to instantiate the exponential mechanism, we give an alternative
inner product approximation protocol that achieves differential privacy.
Using this approximation, we can build a protocol that is able to sample
from exactly the product distribution while additionally leaking a value
$\mathsf{leak}$ that is differentially-private and thus does not violate
the DP guarantee of the exponential mechanism.  
%As discussed earlier,
%the missing piece needed to instantiate the exponential mechanism is a
%product sampling protocol with differentially-private leakage.  
We build such a protocol based on the approximate inner product using
the JLT given in Section~\ref{sec:jlt}.

\paragraph{Approximating the inner product differentially privately.}
We now describe a mechanism, executed by a trusted curator, to
approximate the inner product on inputs $\vec{w}_1$ and $\vec{w}_2$ with
${w}_{b,i} = \frac{{v}_{b,i}}{||\vec{v}_b||_1}$ for $b \in \{1,2\}$ and
$i \in [n]$ as described above.  This mechanism is essentially the {\sf
approxIP} algorithm described in Section~\ref{sec:jlt} with noise added
in the exponent to guarantee differential privacy.

\BEN
\item[] \hspace{-2em} $\dIP(\w_1, \w_2)$: 
 
\item Choose $k \times n$ matrix $\vec{M}$ such that each entry
$M_{i,j}$ chosen from an independent Gaussian distribution of mean $0$
and variance $1$. The dimension $k$ (with $k \ll n$) is determined
appropriately according to Lemma~\ref{lem:jlt_k}.

\item Choose a value $x$ from the Laplace distribution ${\sf
Lap}(1/(\Delta u \cdot |D|))$. 

\item Output $e^x \cdot \langle \vec{M}\vec{w}_1, \vec{M}\vec{w}_2
\rangle$. 
\EEN

\paragraph{Public sampling of $\vec M$.}
Contrary to Section~\ref{sec:rounds} where $\vec M$ was sampled inside
the FHE, here, the matrix $\vec M$ can be publicly sampled (e.g.,
through a commonly chosen random PRG seed), since DP is achieved through
adding a Laplace noise. 

\paragraph{Differential privacy.}
We say that two datasets $D$ and $D'$ are \emph{neighboring} if they
differ in exactly one row.  

Recall that, as in the application described in
Section~\ref{sec:example}, the parties' inputs to the product sampling
are of the form $\vec{w}_b$ where ${w}_{b,i} \propto e^{-\epsilon \cdot
\frac{f^{c_i}_{\mathsf{loss}}(D_b)}{2 \Delta u}}$ for some additive loss functions
$f^{c_i}_\mathsf{loss}$. 


\begin{theorem}
If $|D| \ge 40 \cdot (\log n + \secparam)/\epsilon$, the mechanism
$\dIP$ is $\epsilon$-differentially private w.r.t a database $D_1$
(resp., $D_2$) when the loss functions $f^{c_i}_{\mathsf{loss}}(D)$ have
low sensitivity $1/|D|$.  
\end{theorem}

\begin{proof}
We prove differential privacy with respect to $D_1$ (i.e., the adversary
knows $D_2$ and the differential privacy guarantee holds relative to
rows of $D_1$). The reverse case is analogous.

Note that the change of a single row in $D_1$ causes the values ${v}_{1,i}$ (for $i \in [n]$) and the value $\|\vec{v}_1\|_1$ to change by at most a  
multiplicative factor of $\alpha := e^{\pm \epsilon / (2
\Delta u \cdot |D|)}$.
Thus, letting $\vec{M}_\ell$ be the $\ell$-th row of the JLT matrix (which is
fixed and public), for each $\ell \in [k]$, the value $\langle
\vec{M}_\ell,  \vec{w}_b \rangle$ can change by at most a multiplicative
factor of $\alpha^2$. 
%
Therefore, the dot product $\langle \vec{M} \vec{w}_1, \vec{M}
\vec{w}_2\rangle$ can change by at most a multiplicative factor of
$\alpha^2$. 

Ultimately, this means that we have additive sensitivity of $\epsilon/
(\Delta u \cdot |D|)$ in the \emph{exponent}.  To achieve differential
privacy, we thus need to multiply the dot product estimate by $e^x$,
where $x$ is drawn from $\mathsf{Lap}(1/ (\Delta u \cdot |D|))$.  
\qed
\end{proof}

\paragraph{Correctness.}
We briefly analyze the estimation error due to the added noise.  Since
$x$ is drawn from $\mathsf{Lap}(1/(\Delta u \cdot |D|))$, the
probability that $|x| \geq \epsilon$ is at most $e^{-\Delta u \cdot
|D|}$. When $|D| > \Delta u \cdot \secparam /\epsilon$, this probability
is at most $e^{-\secparam}$, which is negligible. 

In other words, when $D$ is a sufficiently large dataset, with
overwhelming probability, the incurred multiplicative error is $e^{|x|}
\le e^{\epsilon} < 1+2\epsilon$. Thus, differential privacy adds at most
a $1 \pm 2\epsilon$ multiplicative error on top of the error of the
approximation algorithm.

\paragraph{Removing the curator.}
We described the inner product approximation protocol as being run by a
trusted curator.  As is standard, we can replace this curator with a
secure 2-PC evaluating the mechanism to achieve computational DP. 

\subsubsection{Instantiating the Exponential Mechanism}
We now have all the necessary pieces to instantiate a sublinear
communication protocol to evaluate the exponential mechanism for a
database $D$ held jointly by two parties.

\paragraph{The two-party exponential mechanism protocol.}
We now describe the distributed exponential mechanism where $P_b$ has
input $D_b$ and the loss functions have low sensitivity.  This protocol
is in the $\F_\oslone$-hybrid model.

\floatname{algorithm}{Protocol}\label{prot:prodsampdp}
\begin{algorithm}[htbp]
\textbf{Inputs:} Party $P_b$ has input $D_b$   

\begin{enumerate}
    \item $P_b$ computes $\vec{v}_b = (v_{b, 1}, \ldots, v_{b,n})$ such
    that ${v}_{b,i} = e^{-\epsilon \cdot
    \frac{f^{c_i}_{\mathsf{loss}}(D^b)}{2 \Delta u}}$, where $\Delta u =
    \frac{\e}{40(\log n + \secparam)}$. 
    Let $\vec{w}_{b}=\frac{\vec{v}_{b}}{||\vec{v}_b||_1}$.
    
    \item Invoke the $\F_{2PC}$ ideal functionality to evaluate $\eta = {\dIP}(\w_1,
    \w_2)$.  Note that $\eta$ is an $(1+2\epsilon)$-approximation of the inner product.
    
    \item The parties execute the following steps $m=\frac{(1+2\e) \cdot
    \omega(\log{\secparam})}{\eta}$ times in parallel.
    
    \BEN \item Invoke the $\F_\oslone$ ideal functionality with $P_1$ as the sender with input
    $\vec{w}_1$ and $P_2$ as the receiver.  Let $i_{1,1}^j$ and
    $i_{1,2}^j$ be the output of the $j$th execution to $P_1$ and $P_2$
    respectively. 
    
    \item Invoke the $\F_\oslone$ ideal functionality with $P_2$ as the sender with input
    $\vec{w}_2$ and $P_1$ as the receiver.  Let $i_{2,1}^j$ and
    $i_{2,2}^j$ be the output of the $j$th execution to $P_1$ and $P_2$
    respectively.
    
    \EEN
    
    \item Invoke the $\F_{2PC}$ ideal functionality for the following circuit: 
      \BEN
      \item[] \hspace{-2em} Input: $(i_{1,1}^j, i_{2,1}^j, i_{1,2}^j, i_{2,2}^j)$ for $j = 1, \ldots, m$. 
      \item Let $i_1^j = i_{1,1}^j \xor i_{1,2}^j$, $i_2^j = i_{2,1}^j \xor i_{2,2}^j$. 
      \item Find the smallest $j$ such that $i_1^j$ equals $i_2^j$, and
      output $i_1^j$ to both $P_1$ and $P_2$.  If no such $j$ exists,
      output $\mathsf{abort}$.  \EEN

\end{enumerate}
\textbf{Output:} Both parties output the sampled index $i$ or $\mathsf{abort}$. 

\caption{Exponential Mechanism Protocol ($\Pi_{\mathsf{EM}}$) in the $\{\F_{2PC}, \F_\oslone \}$-hybrid model}
\end{algorithm}


\paragraph{Security.}
We will prove the following theorem.

\begin{theorem}\label{thm:prodDP}
If $|D| \ge 40 \cdot \secparam (\log n + \secparam)/\epsilon^2$,
protocol $\Pi_{\mathsf{EM}}$ is
$(2\epsilon,\negl(\secparam))$-DP.
\end{theorem}

It is easy to see that this protocol runs an enough number of product
samplings in parallel so that it does not output $\mathsf{abort}$,
except with negligible probability  (see Section~\ref{sec:const}).
Therefore, for the proof, we assume that the protocol does not output
$\mathsf{abort}$.  

%
%\iffalse
%\begin{claim}\label{lem:prodAbort}
%For any input vectors $\vec{w}_1$, $\vec{w}_2$ such that $\ip{\vec{w}_1}{\vec{w}_2} > 0$\anote{This should probably be the same promise we used before}, the probability that Protocol~\ref{prot:prodsampdp} aborts is negligible in $\secparam$.
%\end{claim}
%\begin{proof}
%We first consider the probability that there is no colliding indices in the first $r-1$ rounds
%\begin{align*}
%\Pr[\mbox{no collision in first }r-1 \mbox{ rounds}] = (1- \ip{\vec{w}_1}{\vec{w}_2})^{r-1}
%\end{align*}
%
%\noindent
%Since $aIP$ is an $(\eps,\delta)$-approximation\anote{Need to get the $\epsilon$ correct} of $\ip{\vec{w}_1}{\vec{w}_2}$, we know that with all but negligible probability $(1-\eps)\ip{\vec{w}_1}{\vec{w}_2} \le aIP \le (1+\eps)\ip{\vec{w}_1}{\vec{w}_2}$.
%
%\medskip\noindent
%Letting $x=1/\ip{\vec{w}_1}{\vec{w}_2}$ and using the inequality $(1-1/x)^x \le e^{-1}$, we see that after $\omega(x\log{\secparam})$ rounds $\Pr[\mathsf{abort}] \le e^{-\omega(\log{\secparam})} \le \negl(\secparam)$.  Thus, given $IP$, this guarantee holds after at most $\frac{IP}{1-\eps}\cdot\omega(\log{\secparam})$ iterations.
%\end{proof}
%\fi
%

\begin{proof}
We first consider where $P_1$ is corrupted by the adversary $A$. 
Let $\view_A(D)$ be the view of the protocol to $A$ (consisting of input, output and transcript)
in the $\{\F_{2PC}, \F_\oslone \}$-hybrid model.
In the $j$-th invocation of $\F_\oslone$,
the outputs $\{i_{1,1}^j,
i_{2,1}^j\}_j, i)$ sent to $A$ by the ideal functionality are uniformly
distributed and independent of $D$.
Thus, WLOG, we assume that 
$A$'s view consists of its input,
transcript $\eta$, and output $i$.
%
Let $\view^{\mathsf{trans}}_A(D)$ (resp., $\view^{\mathsf{out}}_A(D)$) be the transcript (resp., output)
contained in $A$'s view during a random execution of the protocol with input $D$.

For all neighboring database $D$ and $D'$, and for all $\eta$ and $i$, we have: 
\begin{eqnarray*}
&& \Pr[\view^{\mathsf{trans}}_A(D) = \eta, \view^{\mathsf{out}}_A(D) = i] \\
&& = \Pr[\view^{\mathsf{trans}}_A(D) = \eta] \cdot \Pr[\view^{\mathsf{out}}_A(D) = i | \eta]  \\
&& \le e^{\e} \Pr[\view^{\mathsf{trans}}_A(D') = \eta] \cdot \Pr[\view^{\mathsf{out}}_A(D) = i | \eta] \\
&& \le e^{\e} \Pr[\view^{\mathsf{trans}}_A(D') = \eta] \cdot e^{\e} \Pr[\view^{\mathsf{out}}_A(D') = i | \eta] \\
&& = e^{2\e} \Pr[\view^{\mathsf{trans}}_A(D') = \eta, \view^{\mathsf{out}}_A(D') = i],
\end{eqnarray*}

The first inequality holds form the DP of protocol $\dIP$, and the
second inequality holds from the DP of the exponential mechanism.

The case when $P_2$ is corrupted can be proved similarly. 
\qed
\end{proof}



% \paragraph{Performance}
% This protocol executes at most $m=\left(\frac{aIP}{1-\epsilon}\right)\cdot \omega(\log{\secparam})$ executions of the 2-PC comparison protocol in parallel, each requiring $O(\ell)$ communication.  Finally, the 2-PC for selecting the output can be done in communication linear in the number of executions, thus giving a total communication of $O(ml) = o(n)$.  As long as we use constant-round 2-PC protocols to realize all these computations, the resulting protocol also has constant rounds.





